% =================================================
% Учебный лист ученика
% =================================================

\packtitle{Скрытый треугольник и окружность}

\studentnameblock

\section{Основная задача}

\begin{problembox}
В трапеции \(ABCD\) основание \(AD\) в два раза меньше основания \(BC\).
Внутри трапеции взяли точку \(M\) так, что углы \(BAM\) и \(CDM\) прямые.

\begin{enumerate}
  \item[а)] Докажите, что \(BM=CM\).
  \item[б)] Найдите \(\angle ABC\), если \(\angle BCD=64^\circ\), а расстояние
  от точки \(M\) до прямой \(BC\) равно стороне \(AD\).
\end{enumerate}
\end{problembox}
\vspace{2cm}
\trapezoidproblemfigure

\begin{questionsbox}[Первое исследование рисунка]
\begin{enumerate}
  \item Где в геометрии встречается отрезок, параллельный стороне треугольника
  и равный половине этой стороны?
  \item Можно ли включить данную трапецию в более крупный треугольник?
  \item Как могут помочь прямые углы при точках \(A\) и \(D\)?
\end{enumerate}
\end{questionsbox}



\section{Подготовительные задания}

\subsection{Идея 1. Достроить фигуру}

\begin{trainingtask}[A. Скрытый треугольник и средняя линия]
Боковые стороны трапеции продолжили до пересечения в точке \(N\).
Получился треугольник \(BNC\).

\hiddentrianglefigure

Объясните, почему треугольники \(NAD\) и \(NBC\) подобны. Запишите
коэффициент подобия, если \(AD=\frac12BC\).
\answerline[1]

\medskip
На том же рисунке известно
\[
AD\parallel BC,
\qquad
AD=\frac12BC.
\]
Докажите, что \(A\) и \(D\) являются серединами сторон \(BN\) и \(CN\).
Заполните вывод:
\[
AB=\shortanswer{26mm},
\qquad
CD=\shortanswer{26mm}.
\]
\answerline[2]
\end{trainingtask}

\begin{propertybox}[Обратный признак средней линии]
Если отрезок соединяет две стороны треугольника, параллелен третьей стороне
и равен её половине, то его концы являются серединами соответствующих сторон.
\end{propertybox}

\subsection{Идея 2. Увидеть серединный перпендикуляр}

\begin{trainingtask}[C. Два условия, а не одно]
Точка \(A\) — середина отрезка \(BN\), а \(AM\perp BN\).

\perpendicularbisectorfigure

\begin{enumerate}
  \item Как называется прямая \(AM\)?
  \item Какое равенство расстояний отсюда следует?
\end{enumerate}

\[
AM\text{ — }\shortanswer{55mm},
\qquad
MB=\shortanswer{20mm}.
\]
\end{trainingtask}
\begin{trainingtask}[D. Центр окружности]
В треугольнике \(BNC\) прямые \(AM\) и \(DM\) являются серединными
перпендикулярами к сторонам \(BN\) и \(CN\).

\circumcenterfigure

Закончите цепочку:
\[
MB=MN,
\qquad
MC=MN
\quad\Longrightarrow\quad
\boxed{\;MB=\shortanswer{22mm}=\shortanswer{22mm}\;}.
\]
Как называется точка \(M\)? \shortanswer{70mm}
\end{trainingtask}

\newpage
\subsection{Идея 3. Получить угол \(45^\circ\)}

\begin{trainingtask}[E. Перпендикуляр к хорде]
Точка \(M\) равноудалена от точек \(B\) и \(C\). Из неё опустили
перпендикуляр \(ME\) на \(BC\).

\perpendicularmefigure

Почему \(E\) является серединой \(BC\)?
\answerline[2]

Если \(BC=2a\), а \(ME=a\), найдите \(BE\) и \(\angle MBC\).
\answerline[3]
\end{trainingtask}

\trianglebmefigure

\begin{propertybox}[Полезная микроидея]
Если катеты прямоугольного треугольника равны, то его острые углы равны
\(45^\circ\). В нашей задаче
\[
BE=\frac12BC=AD=ME.
\]
\end{propertybox}

\subsection{Идея 4. Перейти от вписанного угла к центральному}

\begin{trainingtask}[F. Одна дуга — два угла]
Точки \(B,C,N\) лежат на окружности с центром \(M\), причём
\(\angle BCN=64^\circ\).

\centralinscribedfigure

\begin{enumerate}
  \item На какую дугу опирается угол \(\angle BCN\)?
  \item Какой центральный угол опирается на ту же дугу?
  \item Найдите этот центральный угол.
\end{enumerate}

\[
\angle BMN=\shortanswer{25mm}.
\]
\end{trainingtask}

\begin{trainingtask}[G. Равнобедренный треугольник \(BMN\)]
Известно, что \(MB=MN\) и \(\angle BMN=128^\circ\).

\trianglebmnfigure

Найдите \(\angle MBN\).
\answerline[2]
\end{trainingtask}

\newpage
\section{Решаем основную задачу}

\subsection{Пункт а. Доказываем равенство \(BM=CM\)}

\begin{stepbox}[Продлите боковые стороны]{1}
Продлите \(AB\) и \(CD\) до пересечения в точке \(N\). Объясните, почему
\(AD\) является средней линией треугольника \(BNC\).
\answerline[3]
\end{stepbox}

\begin{stepbox}[Найдите два серединных перпендикуляра]{2}
Используйте прямые углы \(\angle BAM\) и \(\angle CDM\). Докажите, что
\(AM\) и \(DM\) являются серединными перпендикулярами к сторонам \(BN\)
и \(CN\).
\answerline[4]
\end{stepbox}

\begin{stepbox}[Определите роль точки \(M\)]{3}
Какую точку треугольника \(BNC\) получили? Запишите равенство трёх радиусов
и завершите доказательство пункта а).
\answerline[4]
\end{stepbox}

\Needspace{8\baselineskip}
\begin{importantbox}[Промежуточный результат]
После пункта а) мы знаем больше, чем требовалось:
\[
\boxed{MB=MC=MN}.
\]
Значит, \(M\) — центр окружности, проходящей через вершины \(B,C,N\).
Именно этот дополнительный результат понадобится в пункте б).
\end{importantbox}

\subsection{Пункт б. Находим угол \(ABC\)}

\begin{stepbox}[Опустите перпендикуляр к основанию]{4}
Пусть \(ME\perp BC\). Почему \(E\) — середина \(BC\)? Чему равны
\(BE\) и \(ME\), если \(AD=a\)?
\answerline[3]
\end{stepbox}

\begin{stepbox}[Найдите первый угол]{5}
Докажите, что треугольник \(BME\) прямоугольный равнобедренный. Найдите
\(\angle MBC\).
\answerline[3]
\end{stepbox}

\Needspace{14\baselineskip}
\begin{stepbox}[Найдите второй угол]{6}
Поскольку \(C,D,N\) лежат на одной прямой,
\[
\angle BCN=\angle BCD=64^\circ.
\]
Используйте центральный угол \(\angle BMN\) и равнобедренный треугольник
\(BMN\), чтобы найти \(\angle MBN\).
\answerline[4]
\end{stepbox}

\begin{stepbox}[Сложите найденные части]{7}
Луч \(BM\) проходит внутри угла \(NBC=ABC\).

\anglesumfigure

Запишите последнее равенство решения:
\[
\angle ABC=\angle NBM+\angle MBC
=\shortanswer{15mm}+\shortanswer{15mm}
=\boxed{\shortanswer{18mm}}.
\]
\end{stepbox}

\newpage
\section{Проверяем строгость рассуждения}

\begin{warningbox}[Точки \(N,M,E\) не обязаны лежать на одной прямой]
Из условия \(ME\perp BC\) следует только то, что \(ME\) — единственный
перпендикуляр к \(BC\), проходящий через \(M\). Чтобы утверждать, что
\(N,M,E\) лежат на одной прямой, нужно было бы отдельно доказать
\(MN\perp BC\). В данной задаче это неверно.

\collinearityerrorfigure
\end{warningbox}

\begin{algorithmbox}[Как искать скрытую окружность]
\begin{enumerate}
  \item Достройте исходную фигуру до треугольника.
  \item Найдите середины сторон: часто помогает средняя линия.
  \item Соедините информацию о середине и перпендикуляре.
  \item Распознайте серединные перпендикуляры.
  \item Их пересечение объявляйте центром описанной окружности только после
  проверки обоих условий.
  \item Используйте равенство радиусов и теорему о центральном и вписанном углах.
\end{enumerate}
\end{algorithmbox}

\section{Самостоятельная реконструкция}

\begin{problembox}[Запишите полное решение без подсказок]
Повторно решите основную задачу, используя только следующие ключевые слова:
\[
\begin{gathered}
\text{средняя линия};\qquad
\text{серединный перпендикуляр};\\
\text{центр окружности};\qquad
\text{вписанный угол}.
\end{gathered}
\]
\end{problembox}

\answerline[10]

\section{Обобщение для сильных учеников}

Пусть вместо \(64^\circ\) дано
\[
\angle BCD=\alpha.
\]
Остальные условия задачи не меняются. Выразите \(\angle ABC\) через \(\alpha\).

\Needspace{9\baselineskip}
\begin{hintbox}
Найдите отдельно
\[
\angle MBC
\qquad\text{и}\qquad
\angle MBN.
\]
Первый угол от \(\alpha\) не зависит.
\end{hintbox}

\answerline[5]

\begin{selfcheck}
\begin{itemize}
  \item[\checkbox] Я использовал оба условия для серединного перпендикуляра:
  середину и прямой угол.
  \item[\checkbox] Я объяснил, почему \(M\) — центр описанной окружности.
  \item[\checkbox] Я не предполагал без доказательства, что \(N,M,E\) лежат
  на одной прямой.
  \item[\checkbox] Я указал, на какую общую дугу опираются центральный и
  вписанный углы.
  \item[\checkbox] В ответе пункта б) записано \(71^\circ\).
\end{itemize}
\end{selfcheck}
